% Residencia pública de la genealogía anterior a la publicación parangular.
% Este bloque no sustituye ningún propietario: fija el orden de lectura de
% las dos ramas ya demostradas y de su confluencia tipada.

\section{Bifurcación genealógica de la escala de acción y de su componente adimensional}
\label{sec:sucesor-102-c14-bifurcacion-accion}

La fase intrínseca \(\alpha_{\mathrm{HMT}}\), ya obtenida mediante la
clausura dodecafásica reversible, origina dos ramas internas independientes.
La primera determina la componente adimensional de la acción posterior al
retorno; la segunda determina su década absoluta. Ninguna de ellas presupone
el contraángulo.

En la rama regional, el carácter de grado \(5\) es
\begin{equation}
 H_5(\alpha_{\mathrm{HMT}},\varphi)
 =\frac12\sqrt{\frac{10^3\alpha_{\mathrm{HMT}}}{\varphi}}
  -\alpha_{\mathrm{HMT}}
  +\frac9{16}\alpha_{\mathrm{HMT}}^2
  -\frac59\alpha_{\mathrm{HMT}}^3
  +\frac7{48}\alpha_{\mathrm{HMT}}^4
  -\frac1{54}\alpha_{\mathrm{HMT}}^5 .
 \label{eq:sucesor-102-c14-h5-preangular}
\end{equation}
La completación \(10^3=729+243+27+1\) publica la fase en la coordenada
analítica
\begin{equation}
 x_A=\frac{50\pi}{9}\alpha_{\mathrm{HMT}},
 \qquad
 D(\alpha_{\mathrm{HMT}})=e^{-2x_A}
 =\exp\!\left(-\frac{100\pi}{9}\alpha_{\mathrm{HMT}}\right).
 \label{eq:sucesor-102-c14-determinante-regional}
\end{equation}
Con el selector regional \(\rho_\pi=169\), la recurrencia convergente
\begin{equation}
 \mathscr C_\pi(\alpha)
 =169\alpha^6+
   \frac{D(\alpha)\alpha^7}{1-D(\alpha)\alpha}
 \label{eq:sucesor-102-c14-recurrencia-regional}
\end{equation}
determina directamente
\begin{equation}
 \boxed{\eta_{\mathrm{ret}}
 =H_5-\frac{90}{\pi}\mathscr C_\pi(\alpha_{\mathrm{HMT}})>0.}
 \label{eq:sucesor-102-c14-eta-directa}
\end{equation}
Esta ecuación no contiene \(C^*\): construye la coordenada adimensional de
acción que interviene después en la clausura bicapa.

En la rama determinantal, la forma normal de Smith produce
\begin{equation}
 \Sigma_H=\varphi\alpha_{\mathrm{HMT}}^{16},
 \qquad
 \varepsilon_H=\operatorname{ord}_{10}(\Sigma_H)=-34.
 \label{eq:sucesor-102-c14-decada-accion}
\end{equation}
Las dos ramas confluyen, y sólo entonces, en la magnitud de acción
\begin{equation}
 \boxed{\hbar_{\mathrm{ret}}
 =\eta_{\mathrm{ret}}10^{\varepsilon_H}U_S
 =\eta_{\mathrm{ret}}10^{-34}U_S.}
 \label{eq:sucesor-102-c14-hbar-tipado}
\end{equation}

\section{Publicación parangular y compatibilidad de las cartas global y analítica}
\label{sec:sucesor-102-c14-publicacion-parangular}

Una vez construidas la fase \(\alpha_{\mathrm{HMT}}\), la coordenada
\(\eta_{\mathrm{ret}}\) y la magnitud \(\hbar_{\mathrm{ret}}\), el mapa
parangular publica
\begin{equation}
 \mathcal P(\alpha_{\mathrm{HMT}},\eta_{\mathrm{ret}})
 =\bigl(A_{\deg},C^*_{\deg}\bigr)
 =\bigl([10^3\alpha_{\mathrm{HMT}}]_{360},
        [2(\eta_{\mathrm{ret}}+\alpha_{\mathrm{HMT}})]_{360}\bigr).
 \label{eq:sucesor-102-c14-mapa-parangular}
\end{equation}
El factor \(10^3\) procede de la completación nonádica; el factor \(2\),
de la clausura bidireccional de grado \(2\); y \(360\) es la carta global
que hace simultáneamente visibles
\(12\cdot30=4\cdot90=3\cdot120=360\).
La carta analítica aplica una sola vez
\begin{equation}
 \kappa_{360}(a,c)=\frac{\pi}{180}(a,c),
 \qquad
 A_{\mathrm{rad}}=\frac{50\pi}{9}\alpha_{\mathrm{HMT}},
 \qquad
 C^*_{\mathrm{rad}}
 =\frac{\pi}{90}(\eta_{\mathrm{ret}}+\alpha_{\mathrm{HMT}}).
 \label{eq:sucesor-102-c14-cartas}
\end{equation}
Por tanto,
\(C^*_{\deg}=2(\eta_{\mathrm{ret}}+\alpha_{\mathrm{HMT}})\) es el
teorema directo, mientras que
\(\eta_{\mathrm{ret}}=C^*_{\deg}/2-\alpha_{\mathrm{HMT}}\) constituye
exclusivamente su verificación inversa. La acción satisface
\begin{equation}
 S(\theta)=\hbar_{\mathrm{ret}}\theta_{\mathrm{rad}}
 =h_{\mathrm{ret}}\frac{\theta_{\deg}}{360},
 \qquad h_{\mathrm{ret}}=2\pi\hbar_{\mathrm{ret}}.
 \label{eq:sucesor-102-c14-accion-cartas}
\end{equation}
Los canales \(q_\pm\) evalúan posteriormente los cardinales de incidencia
\(90/120\), ya generados en el Topological Phase Kernel; no constituyen su
origen.
